\documentclass[12pt]{report}
% Definiciones basicas de LaTeX para la memoria
%
% No modificar las lineas a continuacion a menos
% de estar muy seguro de lo que se esta haciendo!!
%
% Jaime Cisternas, julio 2004, agosto 2010
% jecisternas@miuandes.cl
% Sebastian Cea, noviembre 2021
% scea@miuandes.cl 

%%%%%%%%%%%%%%%%%%%%%%%%%%%%%%%%%%%%%%%%%%%%%%%%%%%%%%%%%%%%%%%%%%%%%
% Libreria computin
%\usepackage{algorithm}
%\usepackage{algorithmic}
\usepackage[left=3cm,top=2.5cm,right=3cm,bottom=2.5cm]{geometry}

% Encoding

\usepackage[utf8]{inputenc}

% Carga librerias utiles con simbolos y estructuras
\usepackage{amsfonts}
\usepackage{amssymb}
\usepackage{amsmath}   
\usepackage{amsthm}
\usepackage{latexsym}
\usepackage{cancel}
\usepackage{physics}
\usepackage{bbm}

% Librerías adicionales para diagramas tikz

\usepackage{tikz}
\usepackage{mathdots}
\usepackage{yhmath}
\usepackage{color}
\usepackage{siunitx}
\usepackage{multirow}
\usepackage{gensymb}
\usepackage{tabularx}
\usepackage{extarrows}
\usepackage{booktabs}
\usetikzlibrary{fadings}
\usetikzlibrary{patterns}
\usetikzlibrary{shadows.blur}
\usetikzlibrary{shapes}

% Libreria para ajustar espaciado entre lineas
\usepackage[doublespacing]{setspace}

% Libreria con estilo APA para bibliografia
\usepackage{core/newapa}

% Si usted está utilizando un teclado en español, puede ser útil
% usar los paquetes a continuacion. De otra manera tendrá que generar las
% tildes y la letra ñ de forma especial.
% Estos caracteres no pueden ser utilizados en etiquetas ni tampoco en modo matemático
\usepackage[spanish, activeacute]{babel}
\unaccentedoperators


% Declaraciones especificas de pdfLaTeX
\usepackage[pdftex,
        colorlinks=false,         % true or false (for final version)
        urlcolor=rltblue,         % \href{...}{...} external (URL)
	    anchorcolor=rltbrightblue,
        filecolor=weben,          % \href*{...} local file
        linkcolor=webred,         % \ref{...} and \pageref{...}
        menucolor=webdarkblue,
        citecolor=webgreen,
        pdftitle={},              % INSERT YOUR TITLE HERE
        pdfauthor={},             % INSERT YOUR NAME HERE
        pdfsubject={},
        pdfkeywords={},
        pdfpagemode=None,
        bookmarksopen=true,
	plainpages=false]{hyperref}
%\usepackage[pdftex]{graphicx}

\usepackage{graphicx}


\pdfcompresslevel=9
\pdfadjustspacing=1
% Definicion de colores
\usepackage{color}
\definecolor{rltbrightred}{rgb}{1,0,0}
\definecolor{rltred}{rgb}{0.75,0,0}
\definecolor{rltdarkred}{rgb}{0.5,0,0}
\definecolor{rltbrightgreen}{rgb}{0,0.75,0}
\definecolor{rltgreen}{rgb}{0,0.5,0}
\definecolor{rltdarkgreen}{rgb}{0,0,0.25}
\definecolor{rltbrightblue}{rgb}{0,0,1}
\definecolor{rltblue}{rgb}{0,0,0.75}
\definecolor{rltdarkblue}{rgb}{0,0,0.5}
\definecolor{webred}{rgb}{0.5,.25,0}
\definecolor{webblue}{rgb}{0,0,0.75}
\definecolor{webgreen}{rgb}{0,0.5,0}
\definecolor{webdarkblue}{rgb}{0,0,0.5}
\definecolor{webbrightgreen}{rgb}{0,0.75,0}
% fin de declaraciones especificas pdfLaTeX

% tamanio de pagina
%\setlength{\oddsidemargin}{.5in}
%\setlength{\evensidemargin}{.0in} % en caso de usar opcion 'twoside'
%\setlength{\textwidth}{6in}
%\setlength{\topmargin}{-.5in}
%\setlength{\textheight}{9in}


% Las siguientes definiciones pueden ser usadas en
% memorias mas matematicas
\newtheorem{theorem}{Teorema}[section]
\newtheorem{lemma}[theorem]{Lema}
\newtheorem{corollary}[theorem]{Corolario}
\newtheorem{proposition}[theorem]{Proposición}
\newtheorem{definition}[theorem]{Definición}
\newtheorem{claim}{Afirmación}
\newtheorem{conjecture}[theorem]{Conjetura}
\newtheorem{observation}[theorem]{Observación}
\newtheorem{problem}[theorem]{Problema}

% Definicion de un ambiente para algoritmos en pseudo-codigo.
% Esta definicion puede ser mejorada.
\newtheorem{algorithm}{Algoritmo}[section]
\newcommand{\tab}{\hspace*{0.5 cm}}
% Palablas claves que deben aparecer en negrita
\newcommand{\bfWHILE}{\textbf{while~}}
\newcommand{\bfDO}{\textbf{do~}}
\newcommand{\bfEND}{\textbf{end~}}
\newcommand{\bfIF}{\textbf{if~}}
\newcommand{\bfTHEN}{\textbf{then~}}
\newcommand{\bfELSE}{\textbf{else~}}
\newcommand{\bfFOR}{\textbf{for~}}

% Nombres fijos de Latex pueden ser cambiados al Espaniol
\renewcommand{\contentsname}{Índice General}
\renewcommand{\chaptername}{Capítulo}
\renewcommand{\appendixname}{Anexo}
%\renewcommand{\bibname}{Bibliografía}
\renewcommand{\spanishbibname}{Referencias}
%\renewcommand{\figurename}{Ilustración}
\renewcommand{\figurename}{Figura}
\renewcommand{\tablename}{Tabla}
\renewcommand{\indexname}{Índice}
\renewcommand{\partname}{Parte}
\renewcommand{\listfigurename}{Lista de Ilustraciones}
\renewcommand{\listtablename}{Lista de Tablas}

% Para las enumeraciones usamos primero a,b,c,... y despues i, ii, iii,...
\renewcommand{\labelenumi}{\alph{enumi})}
\renewcommand{\labelenumii}{\roman{enumii})}
\renewcommand{\labelenumiii}{-}
\renewcommand{\labelenumiv}{-}

% Para ser consistentes en el uso de abreviaciones!
%\newcommand{\Chp}{Cap\'{\i}tulo}
%\newcommand{\Chps}{Cap\'{\i}tulos}
%\newcommand{\Sec}{Sec.} % \S
%\newcommand{\Secs}{Secs.} % \S\S
%\newcommand{\SSec}{Subsec.}
%\newcommand{\Fig}{Fig.}
%\newcommand{\Figs}{Figs.}
%\newcommand{\Eqn}{Ecn.}
%\newcommand{\Eqns}{Ecns.}

% Muestra en pantalla los nombres de los archivos usados
\listfiles

% Previene una cierta senial de ``warning: contentsline with no destination''
\newcounter{dummy}

\setlength{\parindent}{0cm}
\usepackage{float}
\usepackage{subfigure}
\usepackage{array}
\newcolumntype{E}{>{$}c<{$}}
\usepackage{longtable}
\setcounter{MaxMatrixCols}{40}
\usepackage{bm}
% Libreria glosarios
% Ver https://es.overleaf.com/learn/latex/Glossaries
\usepackage[acronym,toc]{glossaries}
\makeglossaries

\newglossaryentry{latex}
{
        name=latex,
        description={Is a mark up language specially suited for 
scientific documents}
}

\newglossaryentry{maths}
{
        name=mathematics,
        description={Mathematics is what mathematicians do}
}

\newglossaryentry{formula}
{
        name=formula,
        description={A mathematical expression}
}

\newacronym{gcd}{GCD}{Greatest Common Divisor}

\newacronym{lcm}{LCM}{Least Common Multiple}




%%%%%%%%%%%%%%%%%%%%%%%%%%%%%%%%%%%%%%%%%%%%%%%%%%%%%%%%%%
\newcommand{\nombreautor}{Lucas Moyano}
\newcommand{\mes}{Mes}
\newcommand{\anio}{Año}
\newcommand{\titulo}{Sobre el rol de la inatención racional en inversiones con criterios ASG ante cambios regulatorios}
\newcommand{\nombreprofuno}{Sebastian Cea}
\newcommand{\nombreprofdos}{Nombre Profe 2}
\newcommand{\nombreproftres}{Nombre Profesor Invitado}
\newcommand{\codigo}{ING-IN-XXX/XX}

%%%%%%%%%%Aqui se escribe el resumen%%%%%%%%%%%%%%%%%%%%%

\newcommand{\dedicatoria}{
En esta sección se deben incluir, de ser necesario, los agradecimientos personales a familiares, colegas, etc. (Máximo 1 página)
 }

\newcommand{\resumen}{
En esta sección se debe incluir un resumen de los principales aspectos desarrollados en la tesis, justificación, metodología y principales resultados. (Máximo 1 página). 

\vspace{5cm}

{\color{blue}

El siguiente documento no puede ser la copia del escrito del o los art\'{\i}culos publicados (si corresponde).
Estos deben ser incluidos como Ap\'endices al final de la tesis.

\vspace{1cm}

Los m\'argenes son 2.5 cm superior e inferior, y 3 cm derecha e izquierda, en todo el documento.

\vspace{1cm}

Se debe mantener el formato tipo de letra Times New Roman, tama\~no 12 e interlineado 1.5 l\'{\i}neas.
}
}

\newcommand{\resumenIngles}{
En esta sección se debe incluir un resumen en Inglés de los principales aspectos desarrollados en la tesis, justificación, metodología y principales resultados. (Máximo 1 página).}


\newcommand{\agradecimientos}{
En esta sección, se deben incluir los agradecimientos en primer lugar a las fuentes de financiamiento de las tesis, si las hubiera y posteriormente a los académicos involucrados en el desarrollo de la investigación. (Máximo 1 página). 
}
\renewcommand{\baselinestretch}{1.5}

%%%%%%%%%%%%%%%%%%%%%%%%%%%%%%%%%%%%%%%%%%%%%%%%%%%%%%%%%%

% Por favor coloque todas las definiciones de simbolos a continuacion
% y no en los archivos de capitulos. Esto evitara la  existencia de
% multiples definiciones para las mismas palabras.

\newcommand{\mydef}
	{\stackrel{\mathrm{def}}{=}}
\newcommand{\e}
	{\hbox{\large{e}}}
\newcommand{\mi}
	{\hbox{\large{i}}}
\newcommand{\RR}{\mathbb{R}}


%%%%%%%%%%%%%%%%%%%%%%%%%%%%%%%%%%%%%%%%%%%%%%%%%%%%%%%%%%

\begin{document}
\label{start}

% Genera las paginas de titulo, copyright, resumen,
% agradecimientos, dedicatoria, indices,...

% No es necesario modificar

% Paginas de titulo, derechos de autor y otras mas.

% Jaime Cisternas, julio 2004, agosto 2010

% declaracion especifica de pdfLaTeX
% permite incluir figuras con \includegraphics[...]{figurename}
\DeclareGraphicsExtensions{.jpg,.pdf,.mps,.png}

%%%%%%%%%%%%%%%%%%%%%%%%%%%%%%%%%%%%%%%%%%%%%%%%%%%%%%%%%%%%%%%%

  \pagenumbering{roman}
% \setcounter{page}{1}

% Escribe la pagina de titulo, incluyendo el autor y la facultad

  \thispagestyle{empty}
  \textsc{
  \vspace*{0cm}
  \begin{center}
    \Large
     Universidad de los Andes \\
     Facultad de Ingeniería y Ciencias Aplicadas\\
  \end{center}
     \vspace{1cm}
  \begin{center}
     \includegraphics[angle=0,height=5cm]{logos/Logo-UANDES.png}
  \end{center}
     \vspace{1cm}
  \begin{center}
     \Large
     \titulo
  \end{center}
  \vspace{0.5cm}
  \begin{center}
    \Large
    \nombreautor
  \end{center}
  \vspace{0.5 cm}
  \begin{center}
    Tesis para optar al grado de \\
    MAGÍSTER/DOCTORADO EN CIENCIAS DE LA INGENIERÍA\\
    \vspace{1cm}
    Profesor Guía: \nombreprofuno \\
    Profesor Co-Guía: \nombreprofdos \\
    \vspace{0.5cm}
    \codigo \\
    \vspace{0.5cm}
    Santiago, \mes\ de \anio
  \end{center}
  }

%%%%%%%%%%%%%%%%%%%%%%%%%%%%%%%%%%%%%%%%%%

% Página en blanco
$\,$\\
\newpage
$\,$\\


% Escribe la segunda pagina de titulo para ser firmada por la comision

\cleardoublepage
\thispagestyle{empty}

%\begin{center}
%
%\vspace*{2cm}
%\parbox{10cm}{
%\noindent
%Certifico que he leído esta tesis y que en mi opinión
%su alcance y calidad son completamente adecuados como para ser considerada
%una tesis conducente al título de Ingeniero y grado de Magíster.
%\vspace{1cm}
%
%\hfill
%\begin{tabular}{c}
%\hspace{8cm} \\
%\hline
%\nombreprofuno \\
%(Profesor Guía)
%\end{tabular}
%
%\vspace*{1.5cm}
%
%\noindent
%Certifico que he leído esta tesis y que en mi opinión
%su alcance y calidad son completamente adecuados como para ser considerada
%una tesis conducente al título de Ingeniero y grado de Magíster.
%\vspace{0.75cm}
%
%\hfill
%\begin{tabular}{c}
%\hspace{8cm} \\
%\hline
%\nombreprofdos
%\end{tabular}
%
%\vspace*{1.5cm}
%
%\noindent
%Certifico que he leído esta tesis y que en mi opinión
%su alcance y calidad son completamente adecuados como para ser considerada
%una tesis conducente al título de Ingeniero y grado de Magíster.
%\vspace{0.75cm}
%
%\hfill
%\begin{tabular}{c}
%\hspace{8cm} \\
%\hline
%\nombreproftres
%\end{tabular}
%}
%
%\end{center}

%%%%%%%%%%%%%%%%%%%%%%%%%%%%%%%%%%%%%%%%%%

% Espaciado entre lineas
  \onehalfspacing
  %\doublespacing

% Hace pagina de derechos de autor

%  \cleardoublepage
%  \thispagestyle{empty}
%  %\vspace*{0in}
%  \begin{center}
%    \copyright\ \nombreautor\ \anio \\
%    Todos los derechos reservados.
%  \end{center}

%%%%%%%%%%%%%%%%%%%%%%%%%%%%%%%%%%%%%%%%%%%%%%%%%%%%%%%%%%%%%%%%%%%%%%


% Genera dedicatoria leyendo definicion de agradecimientos

  \cleardoublepage \addcontentsline{toc}{section}{Dedicatoria}
  \begin{center} \Large \textbf{Dedicatoria} \end{center}
  \dedicatoria

% Genera resumen leyendo definicion de resumen

  \cleardoublepage
  \refstepcounter{dummy} \addcontentsline{toc}{section}{Resumen}
  \begin{center} \Large \textbf{Resumen} \end{center}
  \resumen

% Genera resumen leyendo definicion de abstract
  \cleardoublepage
  \refstepcounter{dummy} \addcontentsline{toc}{section}{Abstract}
  \begin{center} \Large \textbf{Abstract} \end{center}
  \resumenIngles



% Genera agradecimientos leyendo definicion de agradecimientos

  \cleardoublepage \addcontentsline{toc}{section}{Agradecimientos}
  \begin{center} \Large \textbf{Agradecimientos} \end{center}
  \agradecimientos

% Produce indice general, listas de figuras y de tablas

  \cleardoublepage
  \tableofcontents
  \cleardoublepage
  \listoffigures
  \cleardoublepage
  \listoftables
  \cleardoublepage

  \normalsize
  \pagenumbering{arabic}

%%%%%%%%%%%%%%%%%%%%%%%%%%%%%%%%%%%%%%%%%%%%%%%%%%%%%%%%%%%


%%%%%%%%%%%%%%%%Ingreso de capitulos%%%%%%%%%%%%%%%%%%%%%%

% espaciado similar a Word 1.5
\spacing{1.4}

% cap1.tex

\chapter{Introducción}
\label{ch:intro} % la etiqueta para referencias

%{\color{blue} En esta sección de debe incluir el fundamento general de la motivación de la investigación desarrollada en la tesis. También se debe incluir la Hipótesis de trabajo, Objetivo General y específicos. (Máximo 5 páginas). Usar mismo tamaño, tipo de letra y espacio del este formato).}

%\vspace{0.5cm}

%The \Gls{latex} typesetting markup language is specially suitable 
%for documents that include \gls{maths}. \Glspl{formula} are 
%rendered properly an easily once one gets used to the commands.

%Given a set of numbers, there are elementary methods to compute 
%its \acrlong{gcd}, which is abbreviated \acrshort{gcd}. This 
%process is similar to that used for the \acrfull{lcm}.

%\vspace{0.5cm}
\section{Motivación}

Hoy en día se ha observado un creciente interés por parte de los inversores en los criterios Ambientales, Sociales y de Gobernanza (ASG o ESG por sus siglas en inglés).\footnote{En la motivación del Proyecto FONDECYT Drivers of Sustainable Finance se citan datos de que los activos domiciliados en Estados Unidos que utilizan estas estrategias representan un tercio del total de la inversión bajo gestión profesional en US SIF (2020).} Esto se debe en gran medida al aumento en la conciencia de los inversionistas en temas los cuales hacen referencia a desafíos ambientales y sociales, tales como el cambio climático, la escasez de recursos naturales, distintas prácticas laborales y la ética empresarial, insistiendo de esta manera a las empresas que reporten de manera mas constante y transparente sobre estos temas \footnote{En el inciso 2 de 2022 in review: the evolution of the ESG reporting landscape}. A nivel global, la inversión en ASG ha experimentado un crecimiento significativo en la ultima década, llegando a aumentar hasta 10 veces su tamaño en el mercado. Según estimaciones de Morningstar (Empresa estadounidense de servicios financieros), los activos totales en fondos designados a ESG superaron los 3,9 millones de dolares a finales de septiembre 2021. Este crecimiento viene de la mano con el incremento de información que las empresas están dando a los inversores. Además, el numero de empresas cotizada que emiten informes corporativos ha aumentado, pasando de menos de 20 en la década del 1990 a mas de 10.000 en la actualidad. \footnote{The Rise of International ESG Disclosure Standards} Asimismo, es intuitivo pensar que las empresas que logran seguir estas políticas consiguen traer a largo plazo un impacto significativo tanto en la sostenibilidad y rendimiento financiero. Por lo tanto, los tomadores de decisiones buscan elegir estas empresas que consiguen buenos resultados en estos tres factores, y de esta forma evitan compañías que basen sus ganancias en vicios (alcohol, tabaco y juegos). Este tipo de preocupaciones puede mejorar o facilitar condiciones para empresas que desarrollen productos y estrategias de inversión sostenible. Por ejemplo, empezando por fondos de inversión socialmente responsables, bonos verdes o hasta acciones de empresas enfocadas en energías renovables y buenas conductas empresariales. Pero las preguntas a plantearse son, ¿Realmente la gente les toma el real peso a estos criterios? ¿Realmente están atentos a los cambios de regulación que han tenido que implementar las empresas para que de esta forma guíen su decisión de si incorporan o no una empresa sobre la otra en el portafolio de inversión?
\\

En Chile, la importancia de la agenda ASG está en constante crecimiento entre las empresas y sus directivos\footnote{https://accionempresas.cl/content/uploads/esg-en-chile-2023.pdf}, aunque aún hay un largo camino por recorrer. Esto comenzó con las empresas más grandes, que, debido a su significativa presencia en el mercado de valores, están más expuestas a estos asuntos. Estas empresas líderes en el ecosistema empresarial se espera que desempeñen un papel crucial al incentivar a otras empresas a adoptar un enfoque de finanzas sostenibles. El desafío actual radica en que tanto las empresas como sus partes interesadas comprendan y asimilen estos conceptos. Esto implica considerar tanto los aspectos financieros como los no financieros de la actividad empresarial en una misma ecuación, y buscar nuevas formas de alcanzar sus objetivos. Es importante destacar que en la actualidad, no solo se trata de crear valor para los accionistas, sino también para los diversos grupos de interés. Se trata de reconocer que los estándares ASG van más allá de cumplir con regulaciones y figurar en clasificaciones. La agenda ASG brinda la oportunidad de generar un valor adicional para todos y de capturar una parte más significativa de ese valor para la empresa. Además, en un entorno donde los clientes se han vuelto más demandantes, es de creciente relevancia entrar mas de lleno en los estándares ASG, con tal de que las empresas puedan ganarse la preferencia y atención de sus clientes. La adopción de prácticas empresariales adecuadas y socialmente aceptables puede influir en la disposición de los consumidores a adquirir productos o servicios.
De esta misma manera, el día 12 de Noviembre del año 2021, la Comisión para el Mercado Financiero emitió una nueva normativa de carácter general N°461, la cual incorpora nuevos temas a la Memoria Anual de las empresas supervisadas, abordando sostenibilidad y gobierno corporativo, con el fin de poder guiar a las empresas a las empresas a reportes ASG mas robustecidos.
\\

Teniendo en cuenta lo anterior, se busca aplicar un modelo de inversiones sostenibles el cual tenga un cambio en una regulación ASG propuesta por el gobierno, y ver la relevancia de la inatención a estos criterios en la valorización de los Retornos esperados en sus carteras de inversión. Es por esto que se busca por medio de los integrantes mencionados al inicio, una base de datos y un catastro donde se pueda analizar en que se están fijando hoy en día los inversores a la hora de querer invertir en sus portafolios, y si realmente, están al tanto de las regulaciones que cada vez toman más importancia en el ámbito financiero. Además, de querer crear un programa donde se reúna que información es relevante para los analistas al momento de querer crear su portafolio y si toman en cuenta estos factores ASG.

\vspace{0.5cm}

Actualmente la inversión sostenible sigue siendo un área confusa de las finanzas que suele interpretarse de distintas maneras. En la mayoría de los casos, se traduce en la creación de carteras de inversión que excluyen categorías consideradas problemáticas, como la "desinversion" en empresas productoras de combustibles fósiles, en un aparente esfuerzo por abordar el cambio climático. Lamentablemente, existe una diferencia entre abstenerse de participar en algo que uno no desea respaldar y luchar activamente contra algo que cree que debe detenerse por el bien común de todos. La desinversion, que a menudo suele confundirse con los boicots, no tiene un impacto claro en el mundo real, ya que el hecho de que el 10\% del mercado no compre sus acciones no equivale a que el 10\% de sus clientes dejen de adquirir sus productos. \footnote{Articulo creado por Tariq Fancy "The Secret Diary of a ‘Sustainable Investor’"}

\vspace{0.5cm}

Además cada vez hay mas intereses en el impacto financiero que se esta teniendo con estos criterios ASG. Los estudios se centran en evaluar como los parámetros ASG pueden afectar el rendimiento y la rentabilidad de las inversiones. Se estan llevando a cabo investigaciones para determinar si las empresas con mejores practicas ASG tienen un mejor desempeño financiero a largo plazo y si los inversores que consideran estos criterios pueden obtener mejores resultados. En Chile, cada vez se están profundizando mas los criterios ASG, como se puede ver en Santander donde cuentan ya con fondos mutuos ASG, apoyando el medioambiente, la sociedad y las buenas prácticas. Con esto, podemos obtener distintas fuentes y bases de datos con tal de poder trabajar y así investigar de mejor forma en que se están basando en su toma de decisiones los inversionistas.

\vspace{0.5cm}

De acuerdo a la investigación realizada en el artículo de \citeA{pastor_sustainable_2021}, “Sustainable investing in equilibrium”, donde se modelan las inversiones realizadas en el mercado teniendo en cuenta los factores ASG (Ambiental, sociales, gobernanza) de las empresas, se demostró que, en equilibrio, los activos verdes (que tienen buenos factores ASG) superan el rendimiento esperado de los activos cafés (no con buenos factores ASG) cuando se producen impactos positivos en el factor ASG, el cual refleja cambios en los gustos de los consumidores por productos ecológicos y en los gustos de los inversores por tenencias verdes. También se hace mención de un modelo el cual de alguna forma calcula el gusto de la gente de querer añadir inversiones “verdes” a sus portafolios con tal de que ha futuro se generen retornos esperados futuros, con el plus de además de no dañar el medio ambiente, u otros factores ASG. 

\vspace{0.5cm}



\section{Hipótesis}

Esta tesis enuncia 2 hipótesis de trabajo descritas a continuación:

\begin{enumerate}
    \item Un cambio regulatorio en los estándares de reporte de métricas de cumplimiento relacionadas a ESG tiene efecto en la valorización de activos financieros. 
    \item Dada la complejidad de las métricas ESG, existe un componente en la valorización de activos relacionado a inatención racional.
\end{enumerate}

\section{Objetivos}

El trabajo se estructura en torno a un objetivo general que articula las hipótesis y N objetivos especificos que ...

\subsection{General} 

Evaluar el rol de la inatención racional de los tomadores de decisión ante un cambio en la regulación sobre reporte de métricas ESG en un modelo de inversión sostenible con varios activos en equilibrio.

\subsection{Específicos} 

Tomando en cuenta la revisión de literatura de \cite{pastor_sustainable_2021} y el modelo propuesto, estructuramos los siguientes objetivos específicos.

\begin{enumerate}
    \item Contextualizar en base a la literatura relacionada considerando los temas de: Compromiso de los inversionistas, subrendimiento de activos sostenibles, sobrerendimiento de activos sostenibles, relación entre inversión sostenible y lealtad cliente-empresa, impacto social de la inversión sostenible, y la caracterización de una frontera eficiente.
    \item Identificar resultados del modelo de \cite{pastor_sustainable_2021} que puedan ser validados con datos
    \item Incorporar inatención racional en la toma de decisiones del inversionista del modelo de \cite{pastor_sustainable_2021}
    \item Construir una base de datos para validar el objetivo específico 2 en dos periodos de tiempo (antes y después del cambio en la regulación)
    \item Analizar el grado de validación de las hipótesis de trabajo
\end{enumerate}


\section{Estructura del informe}

En este capítulo \ref{ch:intro} revisamos.... En el próximo Capítulo \ref{ch:litRev} veremos... En el Capítulo \ref{ch:metodo} ...

{\color{blue} (2 a 3 líneas que informe al lector el contenido de cada capítulo) } % Intro
% cap2.tex

\chapter{Revisión Bibliográfica}
\label{ch:litRev} % la etiqueta para referencias

{\color{blue} En esta sección de debe incluir la revisión general y especifica de la investigación
desarrollada en la tesis.
Aparte de referencias antiguas, las citas deben incluir referencias sustanciales de
los 5 últimos años que den cuenta de la situación actual del tema a investigar,
evitar el uso de páginas web o fuentes no formales de información. (Máximo 15 páginas). }

\vspace{0.5cm}

Para citar bibliografía use \citeA{pastor_sustainable_2020} o bien
\cite{pastor_sustainable_2020}, dependiendo de la redacción.

Para referirse a tablas o a figuras use Tabla \ref{tab:nombreTabla} o Figura \ref{fig:nombreFigura},
y para acrónimos como \gls{latex} definirlos en el archivo core/glosario.tex

\section{5 artículos relacionados}%"Comenzar con las 5 referencias del articulo de Pastor" 
%\vspace{0.5cm}
% Agregar referencias sobre tema 1 
%(Hoeoner et al. (2018) encuentran que el compromiso ESG reduce el riesgo a la baja de las empresas, así como su exposición a un factor de riesgo a la baja.)

%Esto es el abstract traducido: Mostramos que la participación en temas ambientales, sociales y de gobernanza puede beneficiar a los accionistas al reducir los riesgos a la baja de las empresas. Encontramos que las reducciones de riesgo (medidas utilizando el valor en riesgo y los momentos parciales más bajos) varían según los tipos de participación y las tasas de éxito. La participación es más efectiva para reducir el riesgo a la baja cuando se abordan temas ambientales (principalmente el cambio climático). Además, los objetivos con grandes reducciones de riesgo a la baja muestran una disminución de los incidentes ambientales después de la participación. Estimamos que el valor en riesgo de los objetivos de participación disminuye en un 9\% de la desviación estándar después de las participaciones exitosas, en comparación con las empresas de control.

%\vspace{0.5cm}

\subsection{Compromiso de los accionistas}

Por otro lado, se ha demostrado que el compromiso en temas ambientales, sociales y de gobernanza benefecia a los inversionistas reduciendo el riesgo a la baja de las empresas, asi como tambien su exposicion a un factor de riesgo a la baja \cite{hoepner_esg_2022}. Para lograr determinar esto, se utilizaron datos exclusivos proporcionados por una gran institucion inversionista con sede en Reino Unido, siendo esta institución uno de los activistas mas influyentes en temas relacionados al ASG (ver bien a cual institucion especificamente se refiere). De acuerdo a los resultados entregados, los inversionistas principalmente se centran en temas de gobernanza, siendo esta un 51\% de la muestra, enfocandose en temas como la regulacion ejecutiva y la estructura del consejo. En segundo lugar, estan los factores ambientales siendo el tema mas referenciado sobre el riesgo climatico. Por ultimo, se encuentra el factor social, siendo la salud y seguridad, las cadenas de suministro y actos legales los principales temas a mencionar.

\vspace{0.5cm}

Debido a que las interacciones ambientales y sociales tienen una menor probabilidad de ocurrencia, es esperable que al incorporarlas en las carteras de inversiones estas reduzcan el riesgo a la baja (a que se refiere con el riesgo a la baja? podria explicarlo mejor). No obstante, esto no es tan evidente con los temas de gobernanza. Esto se debe a que las interacciones sobre temas de gobernanza podrían tener como objetivo aumentar la toma de riesgos si los gestores no diversificados asumen poco riesgo en comparación con lo óptimo para los accionistas diversificados. 

\vspace{0.5cm}
%Agregar referencias sobre tema 2 
\subsection{Subrendimiento de activos verdes respecto de cafés}
%(Hong and Kacperczyk (2009) Los activos verdes no rinden lo suficiente como los 
%activos cafes.)



Hay ciertos estudios que evidencian lo contrario sobre los beneficios que otorgan los activos verdes. Según la investigación realizada en el articulo de \citeA{hong_price_2009}, estos estudiaron en detalle las acciones de empresas que cotizan en bolsa que se dedican a la producción de alcohol, tabaco, y juegos de azar, y que normalmente son conocidos como acciones inmorales, evaluando además los efectos de las normas sociales en los mercados accionarios. 
\vspace{0.5cm}
Hoy en día hay una cierta barrera social o cultural al momento de querer invertir en estas acciones que promueven vicios en las personas, que por obviedad, tienen una puntuación ASG bajísima comparada a otros activos, igualmente hay evidencia que respalda esto gracias a la adopcion ISR (Inversion socialmente responsable) que han adoptado los administradores de instituciones como fondos de pensiones con el pasar del tiempo. Eso si, las instituciones financieras, al no querer invertir en estas acciones, incurren en un costo financiero al abstenerse en incorporarlas en sus carteras de portafolios. Por un lado, esta el costo de no poder diversificar en companias "pecadoras" que cotizan en bolsa. Por otro, estas acciones suelen ser relativamente baratas (bajos indices precio-valor en libros o precio-ganancias) en comparación con otras. Además, estas acciones "pecaminosas" suelen tener menos supervisión por parte de los analistas financieros en comparaciones a otras que no están enfocadas en vicios. Es por esto mismo que estas acciones, tienden a tener rendimientos esperados mas altos que las otras acciones, que además enfrentan un mayor riesgo de litigio gracias a las normas sociales que están presentes hoy en día.

\vspace{0.5cm}

%Agregar referencias sobre tema 3 
\subsection{Activos verdes superan a los cafés, utilizando definiciones alternativas}
%(Algunos estudios encuentran el resultado opuesto, que los activos verdes superan al cafe, utilizando definiciones alternativas de verde y café)
%Podría hablar de Firms perform better if they are better governed, judging by employee satisfaction (Edmans, 2011)
En la realidad y como va avanzando el mundo con respecto a los factores ESG, se ha demostrado lo opuesto a lo anterior mencionado, evidenciando que las empresas verdes superan en rendimiento a largo plazo a los cafés en diversas definiciones y mediciones de estos factores. Tal como lo estudia \citeA{edmans_does_2011}, el cual demuestra que las empresas al tener mejor gobernanza (siendo juzgadas por la satisfacción de sus empleados) suelen rendir mejor que las que no tienen buenas calificaciones en el mismo ámbito. 

De acuerdo a la evaluación del articulo de Edmans donde se evaluaron las mejores 100 empresas para trabajar en Estados Unidos, estas obtuvieron un alfa de cuatro factores de 3,5\% entre los años 1984 y 2009, estando por encima un 2,1\% sobre otras industrias. Además, estas obtuvieron ganancias sorpresas (a que se refiere con sorpresas?) mas positivas y rendimientos de anuncios significativamente mas positivos, siendo un 1,2\% a 1,7\% mas que otras firmas. Lo anterior mencionado se debe a 3 temas fundamentales que se han ido potenciando con el paso de los años; la importancia del capital humano en las empresas, teniendo en cuenta a los empleados como activos esenciales de la organización y no como objetos reemplazables los cuales pueden crear un valor sustancial en la empresa generando nuevas ideas y formando relaciones humanas; lo incapaz que ha sido el mercado para poder entregar toda la información respecto al verdadero valor de los activos intangibles (tal vez profundizar en lo que se refiere intangible), teniendo en cuenta la falta de información de los inversores y otros factores que hacen que la empresa no este realmente bien valorada en el mercado, generando que los inversores no estén conscientes de los verdaderos beneficios de invertir moralmente bien; y el efecto que se esta teniendo al invertir socialmente responsable (SRI en ingles). Con la teoría tradicional de portafolio de Markowitz en 1959 (arreglar citas), la cual indica que poner limitaciones en las decisiones del inversor (tal como sugiere SRI) reduciría los rendimientos, pues una optimización restringida no es mejor que una que no este restringida. Edmans logro relacionar que ciertos filtros  
%Moskowitz en 1972 (arreglar cita), Luck y Pilotte en 1973 y Derwall Guenster, Bauer, y Koedijk en 2005 
de SRI mejoran los rendimientos si se tratan de periodos cortos, siendo la satisfacción del empleado uno de los factores mas colacionados con la rentabilidad de los retornos, dándole así mas valor a las empresas, y por consiguiente superando a empresas cafés.


\vspace{0.5cm}
%Agregar referencias sobre tema 4 
\subsection{Las inversiones socialmente responsables de una empresa aumentan la lealtad de sus clientes}
%(Albuquerque et al. (2019) construyen un modelo en el cual las inversiones socialmente responsables de una empresa aumenten la lealtad de los clientes.)
Lo anterior también es consistente con lo investigado por \citeA{albuquerque_corporate_2019} los cuales lograron construir un modelo en donde se logra reflejar que las inversiones socialmente responsables (ISR), en empresas que invierten en tener una mayor diferenciación en sus productos, permite un aumento en la lealtad de los clientes, dándole así mas poder a las empresas de poder fijar precios, y de esta manera las empresas tienen estadísticamente y económicamente menos riesgo sistemático y mas valor de mercado, para empresas con un alto nivel de ISR. Esto es representado por un análisis con la Q de Tobin, donde se analiza el el efecto de ISR en el valor de la empresa, dando como resultado una relación positiva entre la Q de Tobin y la ISR, siendo mas fuerte con empresas con mayor diferenciación de productos.
\vspace{0.5cm}
%Agregar referencias sobre tema 5 
\subsection{Las inversiones en ASG tienen un impacto social positivo}
%(Oehmke and Opp (2020)La inversión en ESG tiene un impacto social positivo, siento que el paper es un poco alejado a lo de esta investigacion pero igual lo mencione)

De manera significativa, los activos que están bajo gestión en fondos socialmente responsables han ido creciendo con el tiempo, gracias a que los inversores buscan aumentar su asignación en activos con puntuaciones altas de ASG. Esto es por medio de cambios en los gustos que se están realizando en los clientes e inversores, dando como resultado que los activos verdes estén superando a los activos marrones a largo plazo. En consecuencia a esto, el tamaño de la industria en invertir en estos activos ASG, así también como los alfas de los inversores, dependen directamente de que tanto priorizan a las empresas que sigan estas políticas, de tal forma de sacrificar un poco de rendimiento, obteniendo un excedente de inversión al darse el gusto de elegirlas, generando un impacto social positivo. Esto es consistente con lo mencionado por \citeA{202}Oehmke and Opp (2020) en su articulo ``A Theory of Socially Responsible Investment" (ver bien la cita harvard), donde se realiza un estudio caracterizando las condiciones donde un fondo de inversión socialmente responsable conduce a distintas empresas a reducir externalidades, por medio de restricciones de financiamiento y coordinación entre agentes, generando un impacto y producción mas limpia en la inversión. Esto se aleja un poco a lo nuestro, pero enfatiza el impacto social que se genera, por medio de preferencias en tenencias mas verdes.

\vspace{0.5cm}

%Agregar referencias sobre tema 6 
\subsection{La frontera eficiente de ESG}
%(Pedersen (2021) “Responsible investing: The ESG-efficient frontier”)


Con el avance que ha tenido la inversión teniendo en cuenta los factores ASG, han ido existiendo distintas formas de poder valuar los verdaderos beneficios y costos que podría traer, influyendo de esta forma, en la elección de la cartera de activos y sus precios en equilibrio. Es por esto que \citeA{pedersen_responsible_2021}, logro desarrollar una teoría en la cual las puntuaciones ASG de cada activo se dividía en dos roles. Primero en entregar información sobre los principios de la compañía, y segundo, ayudar en las preferencias del inversor. La solución propuesta por Pedersen es una frontera eficiente ASG, que indica el ratio de sharpe mas lógrale según la cartera determinada. Para poder lograr esto se modelaron tres tipos de inversores:  Primero están los que si tienen conciencia ESG (Tipo-A), los cuales maximizan su utilidad, pero tienen un grado de atención a los factores de ASG al momento de tomar decisiones sobre el riesgo y rendimiento utilizado, eligiendo así su cartera con el mayor ratio de sharpe posible. Luego, están los motivados por ASG (Tipo-M), los que utilizan información exclusivamente ASG para formar sus portafolios y buscan con esto formar una cartera a la derecha de cartera tangente donde se encuentra la frontera eficiente ASG, buscando un equilibrio óptimo entre un gran rendimiento del portafolio, bajo riesgo y altos puntaciones de los factores, explicándose en una optimización a las puntaciones elevadas y el ratio de Sharpe. Por ultimo, están los que no tienen conciencia a los factores ASG (Tipo-U) y solo buscan maximizar su utilidad de media-varianza condicional, por lo que pueden elegir una cartera por debajo de la frontera eficiente ASG, pues se basan en menos información sin contar estos criterios.

\vspace{0.5cm}

Tendre que hablar de la regulacion estado unidense que analizaremos aca en la introduccion tambien?
 % Rev. biblio
% cap3.tex

\chapter{Metodología}
\label{ch:metodo} % la etiqueta para referencias

{\color{blue} En esta sección de debe incluir el diseño experimental y metodología general del trabajo que fueron utilizados, su explicación, así como el detalle de los materiales, equipos y metodologías teóricas y experimentales utilizadas. (Máximo 15 páginas). }

\section{Un modelo de varios activos y retornos ESG}

Sigamos el modelo de \citeB{pastor_sustainable_2020}, donde se tienen $N$ firmas con un retorno sobre el retorno del activo libre de riesgo $r_f$ denotado por $\tilde r_n$ para cada $n=1,...,N$. El vector de retornos de todo el modelo se puede escribir de la siguiente forma

$$\tilde r=(\tilde r_1,...,\tilde r_N)=(\mu_1,...,\mu_N)+(\tilde\epsilon_1,...,\tilde\epsilon_N)=\mu+\tilde\epsilon.$$

Esto quiere decir que la firma $n$ tiene un retorno mañana igual a $\tilde r_n=\mu_n+\tilde\epsilon_n$. En otras palabras, en promedio la firma $n$ tiene un retorno $\mu_n$ pero dado que se realizó un estado de la naturaleza que induce un error $\tilde\epsilon_n$, entonces el retorno es la suma de ambos.

Pregunta: Si cada valor de $\tilde\epsilon_n$ posible refleja un estado de la naturaleza y existe una cantidad $S>0$ de estados posibles: ¿Cómo se puede escribir el vector de retornos de los activos $R_n=(R_n(1),...,R_n(S))$?\\

El vector de retornos de los activos $R_n=(R_n(1),...,R_n(S))$ se escribiría:

$$R_n=(\mu_1 + \tilde\epsilon_1,...,\mu_N + \tilde\epsilon_N)$$

Continuar detalle del modelo hasta 

$$w_m=\int_i\omega_iX_idi$$

\newpage
EJEMPLO:

Aunque $M1$ es un modelo fácil de entender, es bastante difícil de manejar con solvers tradicionales, puesto que es un problema muy no lineal. Alternativamente, se puede definir fácilmente una ecuación de Bellman para $M1$, considerando que es un proceso de decisión secuencial [Nemhauser y Wolsey, 2014]. Aún así, esto podría ser dificultoso para resolver con instancias de gran tamaño (incluso de tamaño mediano), debido a la maldición de la dimensionalidad. No obstante, se propone una transformación preliminar de $M1$ a un modelo de Programación Lineal Entero Mixto (MILP) $M2$, el cual es matemáticamente equivalente y que mantiene los mismos parámetros y se agrega un nuevo grupo de variables binarias. Sea $x_{it} = 1$ si se concede el préstamo al consumidor $i \in I$ en el periodo $t \in \{1,...,T\}$, 0 en otro caso.  \\


$$ \begin{array}{r l l}
	\displaystyle\max_{\vb{x}, \tau, \vb{s}} & z(\vb{x}, \tau, \vb{s}) &   									\\
	\text{s.t.}&  s_{t-1} - s_{t} + \gamma_{t} - \displaystyle\sum_{i \in I} a_{it} x_{it} = 0	& \forall t \in \{1,...,T\}		\\
					&  s_{0} = B & 	                                                                                         	\\
					& x_{it} \leq \dfrac{\tau_{t}}{Pod_{it}} & \forall i \in I , \text{  } \forall t \in \{1,...,T\}                 \\
					& x_{it} \in \{0, 1\} & \forall i \in I , \text{  }  \forall t \in \{1,...,T\}                                    \\
					& \tau_{t} \in \left[0,1 \right] , s_{t} \geq 0 & \forall t \in \{1,...,T\}                                            \\
					
\end{array} $$

%\newpage

\[ z(\vb{x}, \tau, \vb{s}) = \begin{cases} z_1(\vb{x}, \tau, \vb{s}) = \sum_{t=1}^{T}\left(- \alpha\tau_t + \sum_{i \in I} a_{it} x_{it}\right) & \text{Si la IF es con fines de lucro} \\ z_2(\vb{x}, \tau, \vb{s}) = \sum_{t=1}^{T} \left(- \alpha\tau_t + \sum_{i \in I} x_{it}\right) & \text{Si la IF es sin fines de lucro} \end{cases} \] \\

$M2$ es un MILP; por lo tanto, se espera que solvers como Gurobi, Cplex o SCIP (entre otros) puedan manejarlo para instancias grandes. En particular, se resolverá usando SCIP en combinación con Python. Para la solución de este problema, se utilizará el conjunto de datos de Bravo et al. [2013], que tiene miles de prestatarios que solicitan préstamos en un tiempo determinado. \\

Sin embargo, $M2$ asume que la cantidad de clientes que llegan a solicitar un préstamo y la montos solicitados son conocidos, luego se propone un nuevo modelo con parámetros estocásticos.

\section{Modelo $M3$: Modelo con el número de clientes como variable aleatoria}

En $M1$ y $M2$ definimos la demanda como un parámetro, indicando el número de prestatarios $|I|$ y los montos solicitados $a_{it}$, pero ahora en este modelo $M3$ se definirán como esta cantidad de clientes como una variable aleatoria y los montos solicitados como un parámetro estocástico. \\

Se asume que en un periodo $t \in \{1,...,T\}$, los escenarios $k \in \{1,...,K\}$ representan rangos de probabilidad de default en orden incremental. Por lo tanto, si $k_1, k_2 \in \{1,...,K\}$ son tales que $k_1 < k_2$, entonces $pod^{k_1}< pod^{k_2}$. El número de clientes que llegan a solicitar un préstamo es ahora $m_{t}^{k}$ y los montos solicitados por cada prestatario del rango $k$ es $a_{t}^{k}$. Notemos, una vez que un prestatario solicita un préstamo, la IF puede estimar sus probabilidades de default. Se pueden aplicar los procesos estándar de estimación de PD definidos por Baesens et al. [2017b] y Thomas [2009].

Sea $x_{t}^{k} = 1$ si se concede préstamos a los consumidores que pertenecen al rango $k \in \{1,...,K\}$ en el periodo $t \in \{1,\cdots,T\}$, 0 en otro caso. 

Luego notemos que $x_{t}^{k}$ y $\tau_t$ son variables de control, las variables de estado $s_t$ y $m_{t}^{t}$ es una variable aleatoria. Adicionalmente, definimos la función de beneficio $g_t(x_t,\tau_t, m_t)$ y de transición de estado $f_t(x_t, s_{t-1},m_{t})$ como:


\[ g_t(x_t,\tau_t, m_t)  = \begin{cases} g_t^1(x_t,\tau_t,m_t) = \sum_{k=1}^{K}\left(- \frac{\alpha}{K}\tau_t +  \mathbb{E}\{a_{it}\} m_{t}^{k} x_{t}^{k}\right) & \text{Si la IF es con fines de lucro} \\ g_t^2(x_t,\tau_t, m_t) = \sum_{k=1}^{K}\left(- \frac{\alpha}{K}\tau_t +  m_{t}^{k} x_{t}^{k}\right) & \text{Si la IF es sin fines de lucro} \end{cases} \] 
 
$$f_t(x_t,s_{t-1},m_t)= s_{t-1} + \gamma_t - \displaystyle\sum_{k=1}^{K}\mathbb{E}\{a_{it}\} m_{t}^{k} x_{t}^{k}$$

Definiendo el problema en términos de maximización del beneficio esperado

$$ \begin{array}{r l l}
	\displaystyle\max_{\vb{x}, \tau} & \mathbb{E} \left\lbrace \displaystyle\sum_{t=1}^{T} g_t(x_t,\tau_t, m_t) \right\rbrace &   									\\
	\text{s.t.}&  s_{t} = f_t(x_t,s_{t-1},m_t)	& \forall t \in \{1,...,T\}		\\
					&  s_{0} = B &                                                                                   	\\
					& \tau_{t} = \displaystyle\sum_{k=1}^{K} x_{t}^{k}({pod_{t}^k - pod_{t}^{k-1}}) & \forall k \in \{1,...,K\}                 \\
					
					& x_{t}^{k-1} \geq x_{t}^{k} & \forall t \in \{1,...,T\}  , \text{  } \forall k \in \{1,...,K\} \\
					
					& x_{t}^k \in \{0, 1\} & \forall t \in \{1,...,T\}  , \text{  } \forall k \in \{1,...,K\} \\                                      
					& \tau_{t} \in \left[0,1 \right] , s_{t} \geq 0 & \forall t \in \{1,...,T\}                \\
					
\end{array} $$




 % Metodología
\chapter{Resultados y Análisis}
\label{ch:resultados} % la etiqueta para referencias

{\color{blue} En esta sección de debe incluir todos los resultados obtenidos, incluyendo figuras, gráficas y tablas. Estos resultados deben ser explicados, analizados y discutidos respecto del estado del arte de la investigación. (Máximo 20 páginas). } % Resultados
\chapter{Conclusiones}
\label{ch:conclu} % la etiqueta para referencias

{\color{blue} En esta sección de debe incluir todas las conclusiones del trabajo que den respuesta a las diferentes preguntas de investigación, hipótesis y objetivos propuestos. 
La sección de conclusiones no incluye análisis y discusiones de resultados realizados en las secciones anteriores (3 a 4 páginas). }

\section{Conclusiones teóricas}
\section{Conclusiones prácticas}
\section{Investigación futura}

 % Conclusiones

\cleardoublepage

\addcontentsline{toc}{chapter}{Nomenclatura}
\chapter*{Nomenclatura}
%\addcontentsline{toc}{chapter}{Nomenclatura}

%\printglossary[type=\acronymtype]
%\printglossary[title=Special Terms, toctitle=Nomenclatura]

GCD: greatest common divisor

% Ingenieros prefieren la bibliografia antes que los anexos

\cleardoublepage

\addcontentsline{toc}{chapter}{Referencias}
%{\color{blue} En esta sección de debe incluir todas las referencias utilizadas en la tesis. Se debe evitar el uso de páginas web o fuentes no formales de información. (Máximo 15 páginas)

%\vspace{0.5cm}

%Sobre el formato de citas, se debe usar el formato APA, esto es: ``Según datos de Fernández (1999), un 110\% de los conductores no alcanza a frenar... En otro estudio del mismo autor (Fernández, 2000) se establece... Otras investigaciones (Fernández y Vergara, 2000) concluyen además... El trabajo de Fernández et al. (2001) confirma que...''

%\vspace{0.5cm}

%Para mayor información consultar: \url{https://uandes.libguides.com/citasyreferencias} 
%}
\bibliographystyle{newapa}
\bibliography{bibliografia}

\appendix % seguido de archivos de apendices

\cleardoublepage

% Escribe 'Anexos' en el Indice General
\addcontentsline{toc}{chapter}{Anexos}


% anexo_a.tex

\chapter{Primer anexo o apéndice}
\label{app:nombreApendice}

{\color{blue} En esta sección de debe incluir todo el material suplementario utilizado en la tesis, así como los artículos publicados (no incluir enviados o aceptados). }


\section{Sección bla bla ba}

TEXTO
    % archivo de apendice

% \input{ruta_al_Archivo/(apendice_b}    % otro archivo de apendice

\label{end}
\end{document}

%%%%%%%%%%%%%%%%%%%%%%%%%%%%%%%%%%%%%%%%%%%%%%%%%%%%%%%%%%%%
